\chapter{\abstractname}

Automated vulnerability patching promises to close the gap between vulnerability discovery and remediation. While Code Property Graphs (CPG) provide rich structural representation for automated discovery, context-aware code patching requires understanding and adherence of the codebase style and logic. 
Some of the current LLM approaches rely on parametric knowledge acquired during training, which previous studies have shown to be insufficient for complex, real-world patching without external structural guidance and retrieved in-context demonstrations\cite{ye2026well, yin2024thinkrepair, gao2023what, shao2026fix}.
However, the effectiveness of retrieval-augmented, agentic LLM pipelines in addressing this limitation, and even to what extent it occurs in modern LLMs, remain underexplored. 

This thesis presents an empirical study of  structure-aware Retrieval-Augmented Generation (RAG) for vulnerability patching. We evaluated an agentic, tool-calling ReAct patch-generation pipeline by comparing unsupervised structural  embeddings against CodeBERT-based semantic embeddings and their hybrid combination.
The architecture was assessed on two main tasks: vulnerability classification via retrieval and patch generation.
Two datasets of different complexity were used to performance evaluation: AutoPatch, a well-curated dataset, and CVEFixes, real-world for different tasks and very rich dataset.
To compare the impact of the LLM backbones, we benchmarked thee models: GPT-5.3-Codex, security specialized model; GPT-4o, a general purpose cheaper model; and an open weights alternative, DeepSeek-V4-Flash.

The results indicate that structure-aware retrieval is a viable but complementary component of automated patching pipelines. Rather than a superior design, its performance depends on the vulnerability class and the downstream model.

